%%%%%%%%%%%%%%%%%%%%%%%%%%%%%%%%%%%%%%%%%%%%%%%%%%%%%%%%%%%%
%% Chapter 7 — Experimental Evaluation and Discussion
%% (18–28 pages)
%% Purpose: test the model and control hypotheses under matched
%% conditions, EXPLAIN the results rather than only reporting them,
%% and keep robustness, privileged-latent and adaptation claims apart.
%% Main output: the complete evidence chain from offline model quality
%% to closed-loop task performance and adaptation.
%%
%% Every experiment in this chapter evaluates exactly three rows:
%%   PPO-ROM            frozen: command ROM              adaptive: none
%%   PPO-ID-FD          frozen: control ID + strict FD   adaptive: none
%%   PPO-ID(AAM)-FD     frozen: control ID + strict FD   adaptive: AAM + bounded ID adapter
%% There are no other policy variants in the primary comparison.
%%%%%%%%%%%%%%%%%%%%%%%%%%%%%%%%%%%%%%%%%%%%%%%%%%%%%%%%%%%%
\chapter{Experimental Evaluation and Discussion}
\label{chap:evaluation}

\section{Hypotheses and Fairness Controls}
\label{sec:eval:hypotheses}
% - Restate RQ1--RQ5 as testable hypotheses H1--H5:
%   H1: SD_z(s,u,z_true) outperforms matched SD_0 and latent controls on
%       held-out extension and velocity rollouts (already supported).
%   H2: SD_z(s,u,z_hat) retains a statistically supported fraction of the
%       true-latent gain on reserved runs.
%   H3: the true-latent ID adapter improves frozen ID-FD, establishing
%       adapter capacity before AAM error is introduced.
%   H4: the AAM-latent adapter improves frozen ID-FD in paired held-out
%       episodes and approaches the true-latent oracle.
%   H5: adaptation improves mismatch recovery without materially
%       degrading the nominal domain or violating any of the eight gates.
% - Define BEFOREHAND what supports or rejects each hypothesis.
% - Fairness controls: paired factor configurations, terrain meshes, path
%   families, seeds, episode lengths and safety stops across all three
%   stacks; matched policy capacity and training budget.
\todo{write}

\section{Model-Level Experiments}
\label{sec:eval:model-level}
% - ID valve prediction and frozen-ID transfer across nominal,
%   hydraulic, mechanical, combined and held-out domains: report the
%   transfer gap AND the failure mode.
% - FD and ROM one-step and rollout prediction; valve-input ablations
%   and persistence baselines.
% - SD_0 vs SD_z(z_true) vs SD_z(z_hat); AAM latent estimation error and
%   temporal stability under fixed and changing factor configurations.
\todo{write}

\section{System Response}
\label{sec:eval:system-response}
% Matched steps, ramps, PRBS and task commands. Per cylinder: rise time,
% settling time, overshoot, steady-state error and gain; extension,
% velocity and pressure response; action/valve saturation; cross-axis
% coupling. This experiment must EXPLAIN why one stack tracks laterally
% or vertically better, not merely report that it does.
\todo{write}

\section{Delay}
\label{sec:eval:delay}
% Estimate command-to-pressure, command-to-velocity, pressure-to-velocity
% and, where meaningful, cylinder-to-tool response, using controlled
% transients, cross-correlation and frequency/phase analysis. Report
% actuator, environment and AAM-update latencies SEPARATELY. Do not label
% a static kinematic mapping as a temporal delay.
\todo{write}

\section{Nonlinearity}
\label{sec:eval:nonlinearity}
% Sweep command magnitude and sign, pressure, extension, velocity,
% payload, COM and arm posture. Quantify extension/retraction asymmetry,
% dead zone, saturation, pressure-dependent gain, load dependence,
% hysteresis and coupling from repeated PAIRED measurements only.
\todo{write}

\section{End-Effector and Mowing Performance}
\label{sec:eval:tracking}
% Same feasible path families and held-out terrain meshes for all three
% stacks. Report lateral and cut-height RMSE, MAE and p95; maximum and
% along-track error; transient settling after path changes; contact
% fraction; head alignment; impacts; terminations; results per family and
% per worst mesh. The eight-gate contract is the MINIMUM acceptance
% criterion, not a substitute for the full distributions.
\todo{write}

\section{Robustness and Generalization}
\label{sec:eval:robustness}
% In-range anchors, held-out combinations and DECLARED out-of-range
% stress tests. Add observation noise, action delay, sensor dropout,
% terrain shift, payload and COM changes. Separate interpolation from
% extrapolation. Report success rate, tracking degradation, saturation,
% safety terminations and adaptation recovery time relative to nominal.
\todo{write}

\section{Adaptation Ablations}
\label{sec:eval:ablations}
% For PPO-ID(AAM)-FD compare: deployable z-hat; privileged true z; zero,
% mean, shuffled and random latents; frozen ID without adapter; adapter
% without AAM history; shorter/longer history; slower latent update; and
% the optional policy conditioned on z-hat (declare the added input --
% it is NOT the same policy architecture).
% The AAM claim requires deployable z-hat to beat frozen PPO-ID-FD AND to
% approach the true-latent oracle on held-out configurations.
\todo{write}

\section{Statistics and Reporting}
\label{sec:eval:metrics}
% Predeclared primary metrics and seeds. Paired episodes, per-run
% summaries, paired bootstrap confidence intervals. Validation choices
% kept strictly separate from the untouched final test set. Report
% failure cases, worst-case meshes and effect sizes, not only pooled
% means or binary gate verdicts. Report inference cost and throughput.
\todo{write}

\section{Integrated Discussion}
\label{sec:eval:discussion}
% - Explain the MECHANISM behind each positive and negative finding:
%   dataset information, model structure, latent estimation quality,
%   action authority, actuator behaviour.
% - Keep architecture effects, robustness effects and adaptation effects
%   separate.
% - Does the policy mask weak actuator dynamics? Use the standing model
%   error, command magnitude/smoothness and model rollouts as evidence.
% - Include the negative results explicitly: ID-FD's contact failure, the
%   tracking cost of UNCONSTRAINED PPO fine-tuning after distillation,
%   earlier latent redundancy, and the need for wider matched excitation.
% - Report the anchored fine-tuning evidence as a matched before/after on
%   the SAME horizon and mesh set: the warm-start reference checkpoint has
%   never been gate-scored, so the "PPO improved on the student" claim is
%   currently supported by training metrics only, not by the gates.
% - Relate the findings back to the research gap of Chapter~\ref{chap:fundamentals}.
\todo{write}

\section{Threats to Validity}
\label{sec:eval:validity}
% - The deployment plant is itself a learned model (command ROM + PhysX),
%   and PPO-ROM's actuator IS that plant model. State plainly that this
%   makes PPO-ROM the reference condition rather than a third competing
%   approximator, and that the ROM alone was calibrated against PhysX
%   (Section~\ref{sec:models:rom-calibration}) while ID-FD was not.
% - That calibration used the PARKED shapes, not the moving-verge task,
%   so it is weakest under cutter-ground contact.
% - Simscape-to-IsaacLab abstraction mismatch; simulation-to-real limits.
% - Finite factor ranges, sparse design (118 runs, 16 reserved
%   configurations), incomplete operating coverage.
% - Policy stochasticity, model-identifiability limits, and the
%   measurement assumptions behind each gate threshold.
\todo{write}
